\begin{problem}[Radical ideals and reduced closed structures]
Let $I\subseteq A$.  Prove the equivalences
\[
I\text{ radical}
\iff
A/I\text{ reduced}
\iff
I=\bigcap_{\mathfrak p\supseteq I}\mathfrak p.
\]
Show also that
\[
(A/I)_{\mathrm{red}}\cong A/\sqrt I
\]
and explain why $I$ and $\sqrt I$ define the same closed subset of $\Spec A$.
\end{problem}
\paragraph{Why this problem matters.}
This is the precise bridge between classical zero sets and scheme-theoretic multiplicity.
\paragraph{Useful ideas before starting.}
Use $\Nil(A/I)=\sqrt I/I$ and the prime-intersection description of radicals.
\paragraph{Guided hints.}
For the last equivalence, apply the nilradical formula to $A/I$ and translate prime ideals
of the quotient back to primes of $A$ containing $I$.
\begin{solution}
The quotient is reduced exactly when its nilradical is zero.  But
\[
\Nil(A/I)=\sqrt I/I,
\]
so this is equivalent to $\sqrt I=I$.  Also
\[
\sqrt I=\bigcap_{\mathfrak p\supseteq I}\mathfrak p,
\]
which gives the prime-intersection criterion.  Finally
\[
(A/I)_{\mathrm{red}}
=(A/I)/(\sqrt I/I)
\cong A/\sqrt I.
\]
Since prime ideals contain $I$ exactly when they contain $\sqrt I$,
\[
V(I)=V(\sqrt I).
\]
\end{solution}
\paragraph{What happened mathematically.}
Radicalization keeps support and removes nilpotent multiplicity.
\paragraph{Extensions.}
Apply the result to $(x^2y^3)$ in $k[x,y]$.
\paragraph{Provenance.}
New synthesis problem consolidating radical-ideal material from Volume V with the geometric
interpretation assigned to VI/10.
